\section{导读：这不是普通视频生成创业访谈}

本节先说明为什么 EP123 值得做成一份平台经济与 AI 产品方法讲义，而不只是记录 ONE2X 和 Mideo 的创业故事。王冠的主线其实有三层：第一，AI 产品经理在大模型时代的角色如何变化；第二，数据为什么决定智能边界，应用公司如何通过产品内生数据找到机会；第三，生成系统会怎样改写推荐系统、内容平台和创作者/平台/消费者之间的权力分配。

这集最有力的句子是“没有中间商赚差价”。这里的中间商不是某个具体公司，而是互联网平台长期扮演的库存分配者：内容由创作者生产，平台负责分发、匹配、排序、抽成和商业化。生成系统的冲击在于，它把“分配”部分内化到“按需生产”里：用户不一定先进入库存池寻找内容，而是把意图交给生产系统，系统直接生成与用户需求匹配的内容。

\begin{importantbox}{本期核心命题}
AI 时代的信息商品会从“分销平台经济”走向“生产系统经济”。内容不再只是单点作品进入流量池，而可能变成由 recipe、context、模型和用户意图共同生成的连续 IP；平台的货币也可能从注意力转向信任。
\end{importantbox}

\begin{knowledgebox}{视觉策略说明}
本视频是固定访谈画面，没有 slides、白板或产品演示。正文使用封面做来源识别，正文图像全部为自制概念图，用来解释数据三阶段、生成系统栈、推荐系统与生成系统的差异、平台权力迁移等抽象机制。
\end{knowledgebox}

\subsection{本章小结}

EP123 的价值在于把 AI 产品、数据飞轮和内容平台权力结构放在一起讨论。它不是只问“视频生成怎么做”，而是问：当生成系统接管内容生产，平台、创作者、消费者和产品经理的角色会怎样变化。

